\section{Q-Learning under Affiliated Valuations}
\label{sec:exp2_results}

Experiment~1 established the factor hierarchy under identical valuations. Under affiliated values, \citet{MilgromWeber1982} predict that the second-price auction should weakly dominate the first-price auction in revenue. Experiment~2 tests whether Q-learning preserves this ordering by generalising the framework to affiliated valuations with the affiliation parameter $\eta \in \{0, 0.5, 1\}$. The design is a $3 \times 2^3 = 24$ mixed-level factorial crossing four structural factors, namely auction format, affiliation strength ($\eta$), number of bidders, and state information, with Q-learning hyperparameters fixed at levels identified in Experiment~1 ($\alpha = 0.1$, $\gamma = 0.95$, $\varepsilon$-greedy exploration, asynchronous updates, zero initialisation). The primary question is whether valuation interdependence, ranging from purely private values ($\eta = 0$) to near-common values ($\eta = 1$), alters the relative performance of first-price and second-price auctions under algorithmic bidding.

\begin{figure}[H]
  \centering
  \includegraphics[width=\textwidth]{figures/e2_trace}
  \caption{Representative learning trajectory for a single trial of Experiment~2 (first-price auction, $\eta = 0.5$, 2~bidders, 10{,}000 episodes). Rolling-mean revenue converges near the BNE prediction.}
  \label{fig:e2_trace}
\end{figure}

\subsection{Revenue}

Number of bidders is the dominant factor for revenue ($|t| = \ExpTwoRevNbidAbsT$), with the low-to-high shift increasing revenue by \ExpTwoRevTopPctEffect\% of the grand mean. Auction type is statistically significant but economically modest ($|t| = \ExpTwoRevAuctionT$, $p \ExpTwoRevAuctionPFmt$). Table~\ref{tab:exp2_ranked_rev} reports the full effect hierarchy; Figure~\ref{fig:e2_main_rev} displays the directional patterns.

\input{tables/exp2_ranked_rev}

\begin{figure}[H]
  \centering
  \includegraphics[width=0.7\textwidth]{figures/e2_main_rev}
  \caption{Experiment~2: Main effects plot for average revenue. First-price auctions reduce revenue across all factor configurations.}
  \label{fig:e2_main_rev}
\end{figure}

The affiliation parameter $\eta$ has no significant effect on any primary outcome, despite spanning the full range from independent private values to near-common values. This null result contrasts with the \citet{MilgromWeber1982} prediction that affiliation should increase revenue under second-price rules. The number of bidders continues to moderate the auction type effect, with the first-price revenue penalty smaller under four bidders than under two. State information also influences revenue outcomes. The interaction plot (Figure~\ref{fig:e2_int_rev}) displays these moderating relationships.

\begin{figure}[H]
  \centering
  \includegraphics[width=0.7\textwidth]{figures/e2_int_rev}
  \caption{Experiment~2: Interaction plot for average revenue. The auction type $\times$ number of bidders and auction type $\times$ state information interactions are among the strongest two-way effects.}
  \label{fig:e2_int_rev}
\end{figure}

Comparing end-state revenue (final 1{,}000 episodes) with lifetime revenue (averaged over all episodes) reveals an asymmetry between auction formats. Both formats generate higher revenue during the learning phase than at convergence, but first-price auctions show a learning-phase premium roughly three times larger than second-price, offsetting the end-state gap.\footnote{First-price auctions show a learning-phase premium of \ExpTwoFPAPremium{} (mean lifetime revenue \ExpTwoAllMeanFPA{} vs.\ end-state \ExpTwoEndMeanFPA), while second-price auctions show a premium of \ExpTwoSPAPremium{} (\ExpTwoAllMeanSPA{} vs.\ \ExpTwoEndMeanSPA).} In the factorial model for lifetime revenue, the number of bidders and state information remain significant, but auction format is no longer significant, rendering the two formats statistically indistinguishable over the full trajectory.\footnote{Lifetime revenue model: $R^2 = \ExpTwoAllRevRsq$; number of bidders $|t| = \ExpTwoAllRevNbidAbsT$, state information $|t| = \ExpTwoAllRevStateAbsT$, auction format $|t| = \ExpTwoAllRevAuctionAbsT$ ($p \ExpTwoAllRevAuctionPFmt$).} The auction format $\times$ number of bidders interaction persists for lifetime revenue. With two bidders, first-price auctions generate higher lifetime revenue, while with four bidders, second-price auctions lead.\footnote{Interaction: $|t| = \ExpTwoAllRevAuctionxNbidAbsT$, $p \ExpTwoAllRevAuctionxNbidPFmt$. Two bidders: FPA \ExpTwoAllFPATwoBid{} vs.\ SPA \ExpTwoAllSPATwoBid; four bidders: SPA \ExpTwoAllSPAFourBid{} vs.\ FPA \ExpTwoAllFPAFourBid.}

\subsection{Price Volatility}

Unlike Experiment~1, where auction type had minimal effect on volatility, first-price auctions now significantly raise price volatility under affiliated valuations. Table~\ref{tab:exp2_ranked_vol} shows auction type among the significant effects for volatility. The main effects plot (Figure~\ref{fig:e2_main_vol}) confirms the directional increase. The number of bidders and affiliation strength also contribute to volatility differences across configurations.

\input{tables/exp2_ranked_vol}

\begin{figure}[H]
  \centering
  \includegraphics[width=0.7\textwidth]{figures/e2_main_vol}
  \caption{Experiment~2: Main effects plot for price volatility.}
  \label{fig:e2_main_vol}
\end{figure}

Model adequacy diagnostics (Appendix~\ref{sec:appendix_robustness}) confirm that the linear model adequately captures the systematic variation.\footnote{$R^2$ values range from \ExpTwoRsqMin{} to \ExpTwoRsqMax{} across \ExpTwoNobs{} observations. LightGBM cross-validated $R^2$ closely matches OLS for revenue (\ExpTwoRevLGBMRsq{} vs.\ \ExpTwoRevRsq). Of the \ExpTwoNRawSig{} effects significant at $p < 0.05$, \ExpTwoNBHSig{} (\ExpTwoBHPct\%) survive Benjamini--Hochberg correction.}

\subsection{Summary}

The number of bidders dominates all other factors ($|t| = \ExpTwoRevNbidAbsT$). Auction format reduces end-state revenue ($|t| = \ExpTwoRevAuctionT$) but not lifetime revenue, as the end-state gap is offset by a larger learning-phase premium under first-price auctions. The affiliation parameter $\eta$ has no significant impact on any outcome. Global sensitivity analysis confirms the dominance of market thickness for revenue ($S_T = 0.22$) and identifies auction format as the largest contributor to winner's curse frequency ($S_T = 0.36$). Cross-method concordance is low ($\rho = 0.33$), reflecting the small design (24 cells) and the difficulty of distinguishing factor rankings with limited data (Appendix~\ref{sec:appendix_sensitivity}). The failure of the affiliation parameter to influence any outcome raises the question of whether more sophisticated algorithms can exploit the informational structure that Q-learning ignores. Section~\ref{sec:exp3_results} replaces tabular Q-learning with contextual bandits.
